\begin{abstract}
\noindent This study examines how the statistical return to case volume depends on the inferential target in epidemic reproduction number estimation. In a branching process with environmental variation, overdispersed offspring and binomial reporting, a conditionally unbiased estimator from $m$ cases has mean squared error $A/m$ about the current value $R_t$, but $v_R+A/m$ about the super-population process mean $\bar{R}$, because between-period transmission variation is irreducible sampling variation for the process target. Unobserved reporting volatility adds a variance component independent of $m$, while error about the realized mean of finite observed periods decays strictly as $A/(Tm)$. In US influenza hospital admissions across four early-growth seasons (1748 state-weeks), an affine loss model gives a positive conditional intercept with an apparent balance ratio of 37.4 admissions per week (95\% CI 19.1 to 71.0), but an intercept-free power law fits the observed range about as well ($p=0.114$; cross-validation does not favor the affine form). Controlled binomial thinning confirms the $1/m$ scaling of the counting term (rate $c=0.41$, 95\% CI 0.31 to 0.53, vs.\ Poisson benchmark 0.45; 0.48 for $q\ge0.2$) while excluding large intercept changes. State fixed effects yield a within-state ratio of 17.3 (95\% CI 3.3 to 51.4), which is not statistically distinguishable from the cross-sectional fit. Data-matched calibration yields an implied environmental variance of 0.040 (95\% CI 0.026 to 0.056) at a 3.6-day generation-interval proxy, falling to 0.013--0.026 at direct 2.2--2.9-day estimates, with 41\%--51\% of the intercept shared by national weeks. These diagnostics validate the counting-noise model while showing that cross-sectional surveillance cannot distinguish an asymptotic floor from power-law decay or identify transmission variance uniquely.

\medskip
\noindent\textbf{Keywords:}\\
Reproduction number\\
Environmental stochasticity\\
Surveillance design\\
Reporting heterogeneity\\
Influenza hospitalisation\\
Inferential targets
\end{abstract}
